\section{Обоснование решения математической постановки задачи}

В предыдущем разделе математическая постановка задачи была сформулирована как
оптимизационная задача максимизации показателя качества
$Q(k,\pi_k)=\frac{k}{m}$, где $k$ --- длина согласованного префикса трассы, а
$\pi_k = (x_0,\dots,x_k)$ --- поведение модели, объясняющее этот префикс.
Ограничения задачи требуют, чтобы последовательность $\pi_k$ одновременно
являлась допустимым поведением системы переходов и была согласована с первыми
$k$ наблюдениями трассы.

Следовательно, задача состоит в нахождении максимально длинного префикса трассы,
который может быть объяснён допустимым поведением модели. В настоящем разделе
предлагается решение этой задачи и доказывается его корректность.

\subsection{Основная идея решения}

Из формулировки задачи непосредственно следует естественная идея её решения.
Поскольку трасса $\tau = (\sigma_1,\sigma_2,\dots,\sigma_m)$ состоит из
конечного числа наблюдений, поиск максимального согласованного префикса можно
выполнять последовательно, по мере увеличения его длины.

Для этого необходимо:
\begin{enumerate}
    \item определить множество состояний, достижимых после согласования первого
    наблюдения $\sigma_1$;
    \item затем определить множество состояний, достижимых после согласования
    первых двух наблюдений $(\sigma_1,\sigma_2)$;
    \item продолжать этот процесс вплоть до последнего наблюдения $\sigma_m$;
    \item найти наибольший индекс $k$, для которого после согласования первых
    $k$ наблюдений остаётся хотя бы одно достижимое состояние модели.
\end{enumerate}

Если после обработки всех $m$ наблюдений множество достижимых состояний
непусто, то трасса полностью согласована с моделью, а показатель качества равен
единице. Если же на некотором шаге множество достижимых состояний становится
пустым, то дальнейшее продолжение уже невозможно, а значит, согласован лишь
некоторый собственный префикс трассы.

Тем самым задача максимизации показателя качества сводится к последовательному
вычислению достижимых множеств состояний.

\subsection{Рекуррентное построение множеств $D_i$}

Для формализации указанной идеи введём последовательность множеств $D_0, D_1,
\dots, D_m \subseteq X$, где каждое множество $D_i$ интерпретируется как
множество всех состояний модели, достижимых после согласования первых $i$
наблюдений трассы.

Поскольку трасса не содержит отдельного наблюдения начального состояния,
исходное множество определяется только множеством начальных состояний модели.

\begin{definition}\label{def:d0}
Начальное множество достижимых согласованных состояний определяется равенством
\begin{equation}\label{f13}
D_0 = X_0.
\end{equation}
\end{definition}

Это определение означает, что до обработки трассы допустимыми считаются все
начальные состояния модели.

Далее, для каждого $i \in \{1,\dots,m\}$ определим множество $D_i$ рекуррентно.

\begin{definition}\label{def:di}
Для каждого $i \in \{1,\dots,m\}$ множество достижимых согласованных состояний
после обработки первых $i$ наблюдений определяется равенством
\begin{equation}\label{f14}
D_i
=
\left\{
y \in X
:
\exists x \in D_{i-1} : C(\sigma_i,x,y)=1
\right\}.
\end{equation}
\end{definition}

Смысл этого определения состоит в следующем. Состояние $y$ включается в
множество $D_i$ тогда и только тогда, когда существует хотя бы одно состояние
$x \in D_{i-1}$, уже достижимое после согласования первых $i-1$ наблюдений,
такое, что переход из $x$ в $y$ согласован с $i$-м наблюдением $\sigma_i$.

Определение \ref{def:di} отражает принцип пошагового распространения
ограничений. Действительно, каждое следующее наблюдение $\sigma_i$ не
рассматривается изолированно, а применяется только к тем состояниям, которые
уже достижимы после обработки предыдущего префикса трассы.

Это означает, что:
\begin{itemize}
    \item множество $D_0$ описывает все возможные исходные положения системы;
    \item множество $D_1$ описывает все состояния, в которые система может
    перейти после шага, согласованного с $\sigma_1$;
    \item множество $D_2$ описывает все состояния, в которые система может
    перейти после двух последовательных шагов, согласованных с
    $(\sigma_1,\sigma_2)$;
    \item и так далее вплоть до множества $D_m$.
\end{itemize}

\subsection{Корректность построения множеств $D_i$}

Для обоснования корректности предложенного решения необходимо доказать, что
множество $D_i$ действительно имеет именно тот смысл, который ему приписан,
а именно: содержит в точности те состояния, которые достижимы после
согласования первых $i$ наблюдений трассы.

\begin{proposition}
Для любого индекса $i \in \{0,\dots,m\}$ и любого состояния $x \in X$
следующие утверждения эквивалентны:
\begin{enumerate}
    \item $x \in D_i$;
    \item существует последовательность состояний $(x_0,x_1,\dots,x_i) \in
X^{i+1}$, для которой выполнены условия
    \[
    x_0 \in X_0,
    \qquad
    x_i = x,
    \qquad
    \forall k \in \{1,\dots,i\} : C(\sigma_k,x_{k-1},x_k)=1.
    \]
\end{enumerate}
\end{proposition}

Иначе говоря, состояние $x$ принадлежит множеству $D_i$ тогда и только тогда,
когда существует путь длины $i$, начинающийся в допустимом начальном состоянии
и согласованный с первыми $i$ наблюдениями трассы, который оканчивается в
состоянии $x$.

\begin{proof}
Доказательство проводится по индукции по $i$.

\textbf{База индукции.} Пусть $i=0$. Тогда по определению \ref{def:d0} $D_0 =
X_0$. Следовательно, для любого $x \in X$ условие $x \in D_0$ эквивалентно
условию $x \in X_0$. Это, в свою очередь, эквивалентно существованию
последовательности длины $1$, состоящей из единственного состояния $(x)$ и
удовлетворяющей условиям утверждения. База доказана.

\textbf{Переход индукции.}
Пусть утверждение верно для некоторого $i-1 \in \{0,\dots,m-1\}$. Докажем его
для индекса $i$.

\medskip

\noindent
\emph{Необходимость.}
Пусть $x \in D_i$. Тогда по определению \ref{def:di}
\[
\exists y \in D_{i-1} : C(\sigma_i,y,x)=1.
\]

По предположению индукции из условия $y \in D_{i-1}$ следует существование
последовательности состояний $(x_1,x_1,\dots,x_{i-1}) \in X^i$, для которой
выполнены условия
\[
x_0 \in X_0,
\qquad
x_{i-1} = y,
\qquad
\forall k \in \{1,\dots,i-1\} : C(\sigma_k,x_{k-1},x_k)=1.
\]
Добавляя к этой последовательности состояние $x$, получаем последовательность
$(x_2,x_1,\dots,x_{i-1},x) \in X^{i+1}$, для которой дополнительно выполнено
\[
C(\sigma_i,x_{i-1},x)=1.
\]
Следовательно, существует последовательность состояний, удовлетворяющая
условиям пункта б.

\medskip

\noindent
\emph{Достаточность.}
Пусть существует последовательность $(x_1,x_1,\dots,x_i) \in X^{i+1}$, для
которой выполнены условия
\[
x_0 \in X_0,
\qquad
x_i = x,
\qquad
\forall k \in \{1,\dots,i\} : C(\sigma_k,x_{k-1},x_k)=1.
\]
Тогда её префикс $(x_1,x_1,\dots,x_{i-1})$ удовлетворяет условиям для индекса
$i-1$, а потому по предположению индукции $x_{i-1} \in D_{i-1}$. Из условия
$C(\sigma_i,x_{i0},x_i)=1$ и определения \ref{def:di} получаем $x_i \in D_i$.
Так как $x_i = x$, то $x \in D_i$.

Тем самым обе импликации доказаны, следовательно, утверждения 1 и 2
эквивалентны. По принципу математической индукции утверждение доказано для
всех $i \in \{0,\dots,m\}$.
\end{proof}

\subsection{Нахождение оптимального значения показателя качества}

Из доказанного утверждения непосредственно следует, что условие
\[
D_i \neq \varnothing
\]
эквивалентно существованию допустимого поведения модели, согласованного с
префиксом трассы длины $i$. Следовательно, длина максимального согласованного
префикса может быть найдена по формуле
\begin{equation}\label{f:sol-kstar}
k^*(\mathcal{M},\tau)
=
\max
\left\{
i \in \{0,\dots,m\} : D_i \neq \varnothing
\right\}.
\end{equation}

Тем самым оптимальное значение показателя качества равно
\begin{equation}\label{f:sol-Qstar}
Q^*(\mathcal{M},\tau)
=
\frac{k^*(\mathcal{M},\tau)}{m}.
\end{equation}

Следовательно, решение оптимизационной задачи полностью сводится к
последовательному вычислению множеств $D_i$ и последующему поиску наибольшего
индекса, для которого соответствующее множество непусто.

\begin{theorem}\label{th:Qstar}
Оптимальное значение показателя качества согласованности трассы с моделью
определяется формулой
\[
Q^*(\mathcal{M},\tau)
=
\frac{1}{m}
\max
\left\{
i \in \{0,\dots,m\} : D_i \neq \varnothing
\right\}.
\]
\end{theorem}

\begin{proof}
По определению показателя качества
\[
Q(\mathcal{M},\tau)=\frac{k^*(\mathcal{M},\tau)}{m},
\]
где $k^*(\mathcal{M},\tau)$ есть длина максимального префикса трассы, для
которого существует допустимое поведение модели, согласованное с этим
префиксом. Согласно предложению \ref{def:di}, существование такого поведения
для префикса длины $i$ эквивалентно непустоте множества $D_i$. Следовательно,
\[
k^*(\mathcal{M},\tau)
=
\max
\left\{
i \in \{0,\dots,m\} : D_i \neq \varnothing
\right\},
\]
откуда непосредственно следует требуемая формула для $Q^*(\mathcal{M},\tau)$.
\end{proof}

\subsection{Критерий существования согласованного поведения}

Из доказанного индуктивного инварианта непосредственно вытекает основной
критерий решения поставленной задачи.

\begin{theorem}\label{th:decision}
Трасса программы $\tau = (\sigma_1,\sigma_2,\dots,\sigma_m)$ согласована с
моделью $\mathcal{M}$ тогда и только тогда, когда
\begin{equation}\label{f15}
D_m \neq \varnothing.
\end{equation}
Эквивалентно,
\begin{equation}\label{f:sol-fullQ}
Q^*(\mathcal{M},\tau)=1.
\end{equation}
\end{theorem}

\begin{proof}
По определению согласованности трассы с моделью требуется выполнение условия
$\mathrm{B}(\mathcal{M}) \cap \mathrm{M}(\tau) \neq \varnothing$. Это означает,
что существует последовательность состояний
\[
(x_0,x_1,\dots,x_m) \in X^{m+1},
\]
для которой выполнены условия
\[
x_0 \in X_0,
\qquad
\forall i \in \{1,\dots,m\} : C(\sigma_i,x_{i-1},x_i)=1.
\]
Согласно утверждению \ref{def:di}, существование такой последовательности
эквивалентно существованию состояния
\[
x_m \in D_m.
\]
Последнее, в свою очередь, эквивалентно условию $D_m \neq \varnothing$. Далее,
равенство $D_m \neq \varnothing$ означает, что максимальный согласованный
префикс имеет длину не меньше $m$. Так как по определению
$k^*(\mathcal{M},\tau) \le m$, отсюда следует $k^*(\mathcal{M},\tau)=m$.
Следовательно,
\[
Q^*(\mathcal{M},\tau)=\frac{k^*(\mathcal{M},\tau)}{m}=1.
\]
Обратная импликация доказывается аналогично.
\end{proof}

Тем самым задача разрешимости оказывается частным случаем оптимизационной
постановки: необходимо определить, достигается ли наилучшее значение
показателя качества.

\subsection{Алгоритм решения}

Из построения множеств $D_i$ непосредственно следует алгоритм решения как
оптимизационной задачи, так и связанной с ней задачи разрешимости. Алгоритм
последовательно вычисляет множества достижимых согласованных состояний,
находит максимальный индекс $k^*$, для которого соответствующее множество
непусто, и на этой основе вычисляет значение показателя качества.

\begin{algorithm}[htp!]
    \SetAlgoLined
    \KwData{система переходов $\mathcal{M} = (X, X_0, R)$; трасса $\tau = (\sigma_1,\dots,\sigma_m) \in \Sigma^m$; предикат согласованности $C : \Sigma \times X \times X \to \{0,1\}$}
    \KwResult{значения $k^*(\mathcal{M},\tau)$ и $Q^*(\mathcal{M},\tau)$; ответ о полной согласованности трассы}

    $D_0 \gets X_0$\;
    $k^* \gets 0$\;

    \ForEach{$i \in \{1,\dots,m\}$}{
        $D_i \gets \left\{
            y \in X
            :
            \exists x \in D_{i-1} : (x,y) \in R \land C(\sigma_i,x,y)=1
        \right\}$\;

        \If{$D_i = \varnothing$}{
            \textbf{прервать цикл}\;
        }

        $k^* \gets i$\;
    }

    $Q^*(\mathcal{M},\tau) \gets \frac{k^*}{m}$\;

    \If{$k^* = m$}{
        \Return{$k^*$, $Q^*(\mathcal{M},\tau)$, $\tau$ полностью согласована с $\mathcal{M}$}\;
    }
    \Return{$k^*$, $Q^*(\mathcal{M},\tau)$, $\tau$ не полностью согласована с $\mathcal{M}$}\;
    \caption{Нахождение максимального согласованного префикса трассы}
    \label{alg:trace-quality}
\end{algorithm}

Данный алгоритм является корректным в силу теорем \ref{th:Qstar} и
\ref{th:decision}. Если множество состояний $X$ конечно, то алгоритм является
конструктивно осуществимым, так как каждое множество $D_i$ вычисляется за
конечное число шагов. В случае бесконечного пространства состояний для его
реализации требуется дополнительный механизм перебора или символьного
представления множеств состояний и отношения переходов, однако сама
математическая редукция задачи при этом не изменяется.

Особый интерес представляет случай, когда $Q^*(\mathcal{M},\tau) < 1$. Тогда
полная трасса не согласована с моделью. Однако это не означает, что не
существует ни одного согласованного префикса трассы. Напротив, существует
максимальный индекс
\[
k^*(\mathcal{M},\tau)
=
\max
\left\{
i \in \{0,\dots,m\} : D_i \neq \varnothing
\right\},
\]
после которого дальнейшее продолжение становится невозможным.

Смысл числа $k^*(\mathcal{M},\tau)$ состоит в том, что первые
$k^*(\mathcal{M},\tau)$ наблюдений трассы ещё могут быть объяснены некоторым
допустимым поведением модели, тогда как наблюдение
$\sigma_{k^*(\mathcal{M},\tau)+1}$ уже не может быть согласовано ни с одним
продолжением. Тем самым значение $k^*(\mathcal{M},\tau)$ указывает на первый
шаг трассы, на котором возникает расхождение между реализацией и моделью, а
значение $Q^*(\mathcal{M},\tau)$ характеризует степень близости трассы к
полному согласованию.

Это наблюдение важно не только теоретически, но и практически, поскольку
позволяет использовать предложенный метод не только для принятия или
отклонения трассы в целом, но и для количественной оценки степени
несоответствия и для локализации шага, на котором оно возникает.
